\documentclass[11pt,a4paper]{article}
\usepackage[T1]{fontenc}
\usepackage[utf8]{inputenc}
\usepackage[provide=*,lithuanian]{babel}
\usepackage{lmodern,amsmath,amssymb,enumitem}
\usepackage[margin=1.5cm]{geometry}
\setlength{\parindent}{0pt}
\setlist[enumerate]{label=\textbf{\arabic*.},leftmargin=0.65cm,itemsep=0.27cm,topsep=0.2cm,parsep=0pt}
\setlength{\abovedisplayskip}{4pt}
\setlength{\belowdisplayskip}{3pt}
\setlength{\abovedisplayshortskip}{2pt}
\setlength{\belowdisplayshortskip}{3pt}
\begin{document}
\begin{center}
{\Large\bfseries Kvadratų skirtumo formulės taikymas}

\medskip
{\large 60 mišrių uždavinių}
\end{center}
\textbf{Užduotis.} Išspręskite uždavinius. Pasirinkite tinkamą būdą
kvadratų skirtumo formulei pritaikyti.
\medskip
\begin{enumerate}
\item Nekeldami skaičių kvadratu ir neatlikdami daugybos, apskaičiuokite:\[999^2-998\cdot1000.\]
\item Įrodykite tapatybę, taikydami kvadratų skirtumo formulę:\[(a^2+b^2)^2-(a^2-b^2)^2=4a^2b^2.\]
\item Panaikinkite vardiklio iracionalumą ir suprastinkite:\[\frac{\sqrt{11}-3}{\sqrt{11}+3}.\]
\item Žinoma, kad \(p^2-q^2=72\) ir \(p+q=-12\). Raskite \(p-q\).
\item Skaidykite dauginamaisiais kiek įmanoma daugiau (naudokite racionaliuosius koeficientus):\[16x^4-81y^4.\]
\item Raskite natūralųjį skaičių \(n>2\), kuriam:\[n^2-(n-2)^2=44.\]
\item Suprastinkite neapskaičiuodami kvadratų:\[\frac{1004^2-996^2}{1002^2-998^2}.\]
\item Nenaudodami skaičiuotuvo, palyginkite:\[\sqrt{399\cdot401}\quad\text{ir}\quad400.\]
\item Išspręskite lygtį neišskleisdami ketvirtojo laipsnio reiškinių:\[(x^2-4x+3)^2=(2x-2)^2.\]
\item Žinoma, kad \(a^2-b^2=\frac32\) ir \(a+b=3\). Raskite \(ab\).
\item Skaidykite dauginamaisiais kiek įmanoma daugiau (naudokite racionaliuosius koeficientus):\[5a^3-20ab^2.\]
\item Panaikinę vardiklių iracionalumą, apskaičiuokite:\[\frac1{\sqrt5+2}+\frac1{\sqrt5-2}.\]
\item Raskite visas teigiamųjų sveikųjų skaičių poras \((x,y)\), kurioms:\[x^2-y^2=24.\]
\item Neskaičiuodami abiejų reikšmių, palyginkite ir raskite jų skirtumą:\[997^2\quad\text{ir}\quad991\cdot1003.\]
\item Įrodykite tapatybę, taikydami kvadratų skirtumo formulę (\(x\ne0\)):\[(x+\frac1x)^2-(x-\frac1x)^2=4.\]
\item Panaikinkite vardiklio iracionalumą ir suprastinkite:\[\frac{6}{\sqrt{13}-\sqrt7}.\]
\item Skaidykite dauginamaisiais kiek įmanoma daugiau (naudokite racionaliuosius koeficientus):\[x^4-10x^2+9.\]
\item Ar egzistuoja sveikieji skaičiai \(x,y\), kuriems \(x^2-y^2=14\)? Atsakymą pagrįskite.
\item Suprastinkite reiškinį ir nurodykite, su kuriomis kintamųjų reikšmėmis pradinis reiškinys turi prasmę:\[\frac{x^2-9}{x^2+6x+9}.\]
\item Panaikinkite vardiklio iracionalumą ir suprastinkite:\[\frac3{\sqrt7-2}.\]
\item Neapskaičiuodami kvadratų, palyginkite:\[\sqrt{801^2-799^2}\quad\text{ir}\quad40\sqrt2.\]
\item Išspręskite lygtį neišskleisdami ketvirtojo laipsnio reiškinių:\[(x^2-3x+1)^2-(x+1)^2=0.\]
\item Žinoma, kad \(x^2-y^2=-45\) ir \(x-y=-3\). Raskite \(x+y\).
\item Užrašykite trupmeną be šaknų vardiklyje:\[\frac{\sqrt7+\sqrt3}{\sqrt7-\sqrt3}.\]
\item Įrodykite tapatybę, taikydami kvadratų skirtumo formulę:\[(a^2-b^2)(c^2-d^2)=(ac+bd)^2-(ad+bc)^2.\]
\item Skaidykite dauginamaisiais kiek įmanoma daugiau (naudokite racionaliuosius koeficientus):\[(a+b)^2-9(a-b)^2.\]
\item Dviejų teigiamųjų sveikųjų skaičių kvadratų skirtumas yra \(15\). Raskite šiuos skaičius, jei jų sandauga mažiausia.
\item Panaikinkite vardiklio iracionalumą ir suprastinkite:\[\frac5{3+\sqrt{14}}.\]
\item Neatlikdami daugybos ir nekeldami kvadratu, nustatykite, kuris skaičius didesnis ir kiek:\[2024\cdot2032\quad\text{ar}\quad2028^2.\]
\item Išspręskite lygtį neišskleisdami ketvirtojo laipsnio reiškinių:\[4(x^2-x)^2-(x^2+x)^2=0.\]
\item Dviejų skaičių skirtumas yra \(\frac32\), o jų kvadratų skirtumas -- \(\frac{15}{4}\). Raskite šių skaičių sumą.
\item Panaikinkite vardiklio iracionalumą ir suprastinkite:\[\frac{\sqrt6}{\sqrt6-\sqrt2}.\]
\item Skaidykite dauginamaisiais kiek įmanoma daugiau (naudokite racionaliuosius koeficientus):\[(x^2+2)^2-9x^2.\]
\item Raskite visas teigiamųjų sveikųjų skaičių poras \((x,y)\), kurioms:\[x^2-y^2=33.\]
\item Neapskaičiuodami kvadratų, raskite reiškinio reikšmę:\[\frac{1003^2-997^2}{2000}.\]
\item Panaikinkite vardiklio iracionalumą ir suprastinkite:\[\frac{\sqrt5-1}{\sqrt5+1}.\]
\item Įrodykite tapatybę, taikydami kvadratų skirtumo formulę:\[(x-y)(x+y)+(y-z)(y+z)=(x-z)(x+z).\]
\item Suprastinkite reiškinį ir nurodykite, su kuriomis kintamųjų reikšmėmis pradinis reiškinys turi prasmę:\[\frac{4a^2-25b^2}{2a^2+5ab}.\]
\item Išspręskite lygtį neišskleisdami ketvirtojo laipsnio reiškinių:\[(x^2+x-3)^2-(2x^2-x-3)^2=0.\]
\item Panaikinkite vardiklio iracionalumą ir suprastinkite:\[\frac1{\sqrt{10}-3}-\frac1{\sqrt{10}+3}.\]
\item Žinoma, kad \(x^2-y^2=48\) ir \(x-y=4\). Raskite \(xy\).
\item Neatlikdami daugybos, nustatykite, kuri sandauga didesnė ir kiek:\[1001\cdot999\quad\text{ar}\quad1002\cdot998.\]
\item Įrodykite tapatybę, taikydami kvadratų skirtumo formulę:\[(a+b)^2-(a-b)^2=4ab.\]
\item Panaikinkite vardiklio iracionalumą ir suprastinkite:\[\frac2{\sqrt7+\sqrt5}+\frac2{\sqrt7-\sqrt5}.\]
\item Žinoma, kad \(a^2-b^2=60\) ir \(a+b=10\). Raskite \(ab\).
\item Suprastinkite reiškinį ir nurodykite, su kuriomis kintamųjų reikšmėmis pradinis reiškinys turi prasmę:\[\frac{x^4-16}{x^2-4}.\]
\item Išspręskite lygtį neišskleisdami ketvirtojo laipsnio reiškinių:\[(2x^2-5)^2-(x^2+1)^2=0.\]
\item Panaikinkite vardiklio iracionalumą ir suprastinkite:\[\frac{\sqrt{18}-\sqrt8}{\sqrt3-1}.\]
\item Neatlikdami triženklių skaičių daugybos, apskaičiuokite:\[1003\cdot997-1001\cdot999.\]
\item Žinoma, kad \(u^2-v^2=-24\) ir \(u+v=6\). Raskite \(uv\).
\item Suprastinkite reiškinį ir nurodykite, su kuriomis kintamųjų reikšmėmis pradinis reiškinys turi prasmę:\[\frac{x^2-1}{x^2-4x+3}\cdot\frac{x-3}{x+1}.\]
\item Panaikinkite vardiklio iracionalumą ir suprastinkite, kai \(a\ge0\):\[\frac1{\sqrt{a+1}-\sqrt a}.\]
\item Dviejų teigiamųjų skaičių skirtumas yra \(3\), o jų kvadratų skirtumas -- \(15\). Raskite jų sandaugą.
\item Suprastinkite reiškinį ir nurodykite, su kuriomis \(x\) reikšmėmis pradinis reiškinys turi prasmę:\[\frac{x^2-4}{x^2+4x+4}:\frac{x-2}{x+3}.\]
\item Suprastinkite taip, kad vardiklyje neliktų šaknų, kai \(a>0\) ir \(b>0\):\[\frac{a-b}{\sqrt a+\sqrt b}.\]
\item Žinoma, kad \(a^2-b^2=84\) ir \(a+b=14\). Raskite \(a-b\).
\item Apskaičiuokite kuo trumpesniu būdu:\[497^2-503^2+6000.\]
\item Panaikinkite vardiklio iracionalumą ir suprastinkite (\(x\in\mathbb R\)):\[\frac4{\sqrt{x^2+4}+x}.\]
\item Žinoma, kad \((2x+1)^2-(y-3)^2=56\) ir \(2x+y-2=14\). Raskite \(2x-y+4\).
\item Panaikinkite vardiklio iracionalumą ir suprastinkite:\[\frac1{\sqrt3+\sqrt2+1}.\]
\end{enumerate}
\clearpage
\section*{Atsakymai}
\begin{enumerate}
\item \(1\).
\item Tapatybė teisinga.
\item \(10-3\sqrt{11}\).
\item \(p-q=-6\).
\item \((2x-3y)(2x+3y)(4x^2+9y^2)\).
\item \(n=12\).
\item \(2\).
\item \(\sqrt{399\cdot401}<400\).
\item \(x_1=1,\quad x_2=5\).
\item \(ab=\dfrac{35}{16}\).
\item \(5a(a-2b)(a+2b)\).
\item \(2\sqrt5\).
\item \((x,y)=(7,5)\) arba \((5,1)\).
\item \(991\cdot1003<997^2\).
\item Tapatybė teisinga.
\item \(\sqrt{13}+\sqrt7\).
\item \((x-1)(x+1)(x-3)(x+3)\).
\item Ne.
\item \(\dfrac{x-3}{x+3}\), kai \(x\ne-3\).
\item \(\sqrt7+2\).
\item \(\sqrt{801^2-799^2}=40\sqrt2\).
\item \(x_1=0,\quad x_2=4\).
\item \(x+y=15\).
\item \(\dfrac{5+\sqrt{21}}2\).
\item Tapatybė teisinga.
\item \(4(2b-a)(2a-b)\).
\item Skaičiai \(4\) ir \(1\); mažiausia sandauga \(4\).
\item \(\sqrt{14}-3\).
\item \(2024\cdot2032<2028^2\).
\item \(x_1=0,\quad x_2=\frac13,\quad x_3=3\).
\item \(\dfrac52\).
\item \(\frac{3+\sqrt3}{2}\).
\item \((x-1)(x-2)(x+1)(x+2)\).
\item \((x,y)=(17,16)\) arba \((7,4)\).
\item \(6\).
\item \(\frac{3-\sqrt5}{2}\).
\item Tapatybė teisinga.
\item \(\dfrac{2a-5b}{a}\), kai \(a\ne0\) ir \(2a+5b\ne0\).
\item \(x_1=-\sqrt2,\quad x_2=0,\quad x_3=\sqrt2,\quad x_4=2\).
\item \(6\).
\item \(xy=32\).
\item \(1002\cdot998<1001\cdot999\).
\item Tapatybė teisinga.
\item \(2\sqrt7\).
\item \(ab=16\).
\item \(x^2+4\), kai \(x\ne\pm2\).
\item \(x_1=-\sqrt6,\quad x_2=-\dfrac{2\sqrt3}{3},\quad x_3=\dfrac{2\sqrt3}{3},\quad x_4=\sqrt6\).
\item \(\frac{\sqrt6+\sqrt2}{2}\).
\item \(-8\).
\item \(uv=5\).
\item \(1\), kai \(x\notin\{-1;1;3\}\).
\item \(\sqrt{a+1}+\sqrt a\), kai \(a\ge0\).
\item \(4\).
\item \(\dfrac{x+3}{x+2}\), kai \(x\notin\{-3;-2;2\}\).
\item \(\sqrt a-\sqrt b\), kai \(a>0\) ir \(b>0\).
\item \(a-b=6\).
\item \(0\).
\item \(\sqrt{x^2+4}-x\), \(x\in\mathbb R\).
\item \(2x-y+4=4\).
\item \(\dfrac{2+\sqrt2-\sqrt6}{4}\).
\end{enumerate}
\end{document}
